% manuscript-sha256: 0dde974e14ff1b54ff6e65705a1194aea38fad7e65bf77b159a07ea1b96bb9a5
% Options for packages loaded elsewhere
\PassOptionsToPackage{unicode}{hyperref}
\PassOptionsToPackage{hyphens}{url}
\PassOptionsToPackage{dvipsnames,svgnames,x11names}{xcolor}
\documentclass[
  11pt,
]{article}
\usepackage{xcolor}
\usepackage[margin=25mm]{geometry}
\usepackage{amsmath,amssymb}
\setcounter{secnumdepth}{-\maxdimen} % remove section numbering
\usepackage{iftex}
\ifPDFTeX
  \usepackage[T1]{fontenc}
  \usepackage[utf8]{inputenc}
  \usepackage{textcomp} % provide euro and other symbols
\else % if luatex or xetex
  \usepackage{unicode-math} % this also loads fontspec
  \defaultfontfeatures{Scale=MatchLowercase}
  \defaultfontfeatures[\rmfamily]{Ligatures=TeX,Scale=1}
\fi
\usepackage{lmodern}
\ifPDFTeX\else
  % xetex/luatex font selection
  \setmainfont[]{Latin Modern Roman}
  \setmonofont[]{Latin Modern Mono}
\fi
% Use upquote if available, for straight quotes in verbatim environments
\IfFileExists{upquote.sty}{\usepackage{upquote}}{}
\IfFileExists{microtype.sty}{% use microtype if available
  \usepackage[]{microtype}
  \UseMicrotypeSet[protrusion]{basicmath} % disable protrusion for tt fonts
}{}
\makeatletter
\@ifundefined{KOMAClassName}{% if non-KOMA class
  \IfFileExists{parskip.sty}{%
    \usepackage{parskip}
  }{% else
    \setlength{\parindent}{0pt}
    \setlength{\parskip}{6pt plus 2pt minus 1pt}}
}{% if KOMA class
  \KOMAoptions{parskip=half}}
\makeatother
\usepackage{longtable,booktabs,array}
\usepackage{calc} % for calculating minipage widths
% Correct order of tables after \paragraph or \subparagraph
\usepackage{etoolbox}
\makeatletter
\patchcmd\longtable{\par}{\if@noskipsec\mbox{}\fi\par}{}{}
\makeatother
% Allow footnotes in longtable head/foot
\IfFileExists{footnotehyper.sty}{\usepackage{footnotehyper}}{\usepackage{footnote}}
\makesavenoteenv{longtable}
\setlength{\emergencystretch}{3em} % prevent overfull lines
\providecommand{\tightlist}{%
  \setlength{\itemsep}{0pt}\setlength{\parskip}{0pt}}
\usepackage{bookmark}
\IfFileExists{xurl.sty}{\usepackage{xurl}}{} % add URL line breaks if available
\urlstyle{same}
\hypersetup{
  colorlinks=true,
  linkcolor={Maroon},
  filecolor={Maroon},
  citecolor={Blue},
  urlcolor={Blue},
  pdfcreator={LaTeX via pandoc}}

\author{}
\date{}

\usepackage{xurl}
\usepackage{newunicodechar}
\usepackage{needspace}
\newunicodechar{∎}{\ensuremath{\square}}
\newunicodechar{→}{\ensuremath{\to}}
\setlength{\emergencystretch}{3em}
\begin{document}

\begin{center}
{\LARGE\bfseries Before the First Mark\par}
Jeremy H. Carroll\par
2026-10-09\par
Version 1.0.0\par
\end{center}

\subsection{Abstract}\label{abstract}

Before interpreting \texttt{\textless{}\textgreater{}}, distinguish
receiving an inscription from admitting an object and distinguish an
external reader's resources from a regime's operations. This paper
combines the source-grounded Pre-Genesis experiments with the Genesis
account of finite admission, observation and response selection. The
experimental record contains three runs, with no recorded model
identifier or generation settings; it supplies examples of failure, not
performance estimates. The mathematical account gives exact criteria for
recovering answers from reports, a two-item separation minimum,
classification of its ties, and a finite closure bound. It also records
the closed Genesis continuation study and its proofs of admission growth
and no return, while preserving the source's unresolved completion
obligations. These results support precise questions about presentations
and question-relative report reconstruction. A separately reported Rust
construction connects primitive admissions with finite incidence
complexes and exact rational homology. The standalone replay in this
package covers the primitive layer; the construction's execution
correspondence and chosen mathematical interpretation have separate
scopes.

\subsection{\texorpdfstring{1. Starting before
\texttt{\textless{}\textgreater{}}}{1. Starting before \textless\textgreater{}}}\label{starting-before}

The first question is what has been supplied and what has been licensed
to happen with it. The author, reader and computer already use a
metalanguage. This paper makes that fact explicit; it does not describe
an observer-free creation from nothing.

\begin{longtable}[]{@{}
  >{\raggedright\arraybackslash}p{(\linewidth - 4\tabcolsep) * \real{0.3333}}
  >{\raggedright\arraybackslash}p{(\linewidth - 4\tabcolsep) * \real{0.3333}}
  >{\raggedright\arraybackslash}p{(\linewidth - 4\tabcolsep) * \real{0.3333}}@{}}
\toprule\noalign{}
\begin{minipage}[b]{\linewidth}\raggedright
Condition
\end{minipage} & \begin{minipage}[b]{\linewidth}\raggedright
Host-level description
\end{minipage} & \begin{minipage}[b]{\linewidth}\raggedright
What has not followed
\end{minipage} \\
\midrule\noalign{}
\endhead
\bottomrule\noalign{}
\endlastfoot
Absent input & No payload was supplied, for example \nolinkurl{None} in a
declared option type. & No object-level empty word or ground has been
produced. \\
Empty word & The word epsilon in a declared free monoid has length zero.
& It need not be accepted by a selected grammar start. \\
Literal \texttt{\textless{}\textgreater{}} & Two ASCII characters,
possibly one opaque token under a stated lexer. & Neither an empty set
nor a singleton carrier is forced. \\
Admitted ground & A named initial object under an explicit regime
definition. & Its carrier, operations and permitted transitions still
need definitions. \\
\end{longtable}

``Before'' here marks the order of questions in the account. It is not a
derivation of physical time. Even writing \nolinkurl{None} presupposes a
host capable of representing absence. Likewise, a grammar production
cannot demonstrate that there is no grammar in the metalanguage in which
that production was written.

Let a declared transition system have state set X and transition
relation T. If T is empty, no state has a T-successor. This follows from
the definition of an empty relation. It does not say that an external
caller cannot provide more data, or that every later rule is prohibited.
External supply and admission are events whose effects must be
specified.

\subsection{2. What the Pre-Genesis trials actually
recorded}\label{what-the-pre-genesis-trials-actually-recorded}

The \href{docs/sources/pre_genesis/readme.txt}{Pre-Genesis strand}
contains the experiment design, Trial 1, Trial 2 and a full-history
replay. Its \href{docs/sources/pre_genesis/02_trial_record.txt}{trial
record} points to the original transcripts. The prompts withheld the
intended comparison theories and asked the model to distinguish external
presentation from internal expressibility. ``No move is forced'' was an
admissible outcome. The next inscription was not prescribed as a
successful rediscovery of an existing calculus.

The record supports the following limited observations:

\begin{longtable}[]{@{}
  >{\raggedright\arraybackslash}p{(\linewidth - 4\tabcolsep) * \real{0.3333}}
  >{\raggedright\arraybackslash}p{(\linewidth - 4\tabcolsep) * \real{0.3333}}
  >{\raggedright\arraybackslash}p{(\linewidth - 4\tabcolsep) * \real{0.3333}}@{}}
\toprule\noalign{}
\begin{minipage}[b]{\linewidth}\raggedright
Event
\end{minipage} & \begin{minipage}[b]{\linewidth}\raggedright
Recorded behavior
\end{minipage} & \begin{minipage}[b]{\linewidth}\raggedright
Evidential boundary
\end{minipage} \\
\midrule\noalign{}
\endhead
\bottomrule\noalign{}
\endlastfoot
Bare marks and punctuation controls & A fenced prompt elicited familiar
interpretations; raw controls elicited different replies. & The runs did
not isolate every causal difference between conditions. \\
No declared continuation rule & All three runs said that a successor was
not forced; some proposed a continuation under an added assumption. & A
proposal and its adoption are different events. \\
\texttt{\textless{}a\textgreater{}} supplied after
\texttt{\textless{}\textgreater{}} & Replies differed between
presentation-relative and inscription-relative readings of identity. &
The identity convention had not been fixed. \\
A separation task about \texttt{\textless{}a\textgreater{}} and
\texttt{\textless{}ab\textgreater{}} & A reply treated observational
indistinguishability as identification; another used a stale account of
the current state. & Lack of a distinguishing operation does not itself
install an object-language equality rule. \\
Truncated audit versus replay & An auditor examining only the local
response missed a historical-state mismatch. & Local consistency alone
did not establish consistency with the adopted history. \\
\end{longtable}

Trial 1 also bundled steps and fenced controls that the design had
intended to send separately and raw. These are recorded design
deviations. The model identifier, settings and dates were not recorded.
There were no host-checker verdicts on the replies and no live PFP
evaluation. The earlier invented ``50 trials'' table is not part of the
evidence used here. No rate, benchmark advantage or theory-rediscovery
claim follows from this record.

The
\href{docs/sources/pre_genesis/04_named_failure_modes.txt}{failure-mode
chapter} names semantic overcompletion, premature relation, false
provenance, lost state, stale certification and bundled extension. These
are retrospective descriptions. Naming a failure does not implement a
detector for it.

\subsection{3. Six operations that must remain
distinct}\label{six-operations-that-must-remain-distinct}

The trials become easier to interpret when the following operations have
separate signatures and records:

\begin{longtable}[]{@{}
  >{\raggedright\arraybackslash}p{(\linewidth - 2\tabcolsep) * \real{0.5000}}
  >{\raggedright\arraybackslash}p{(\linewidth - 2\tabcolsep) * \real{0.5000}}@{}}
\toprule\noalign{}
\begin{minipage}[b]{\linewidth}\raggedright
Operation
\end{minipage} & \begin{minipage}[b]{\linewidth}\raggedright
Meaning in a declared host
\end{minipage} \\
\midrule\noalign{}
\endhead
\bottomrule\noalign{}
\endlastfoot
Mention & Refer to an inscription or possible object in a discussion. \\
Supply & Deliver a payload to an interface. \\
Designate & Identify which presentation is currently under
consideration. \\
Admit & Add a datum, form, instance or rule according to an admission
policy. \\
Retain & Preserve an admitted object or event for subsequent use. \\
Operate & Apply an admitted transformation with stated preconditions. \\
\end{longtable}

A response proposing \texttt{\textless{}a\textgreater{}} does not
establish that \texttt{\textless{}a\textgreater{}} was adopted. A later
task that mentions \texttt{\textless{}ab\textgreater{}} need not erase
\texttt{\textless{}a\textgreater{}}. A checker cannot reconstruct these
facts from punctuation alone; it needs event semantics and enough of the
event record.

The trial method is consequently modest: keep discovery prompts free of
expected solutions, preserve the exact controls, separate steps, record
adoption explicitly, and bind an audit to its actual context. These are
design lessons from the record, not experimentally established universal
remedies.

\subsection{4. Observation-relative
difference}\label{observation-relative-difference}

Let A be a declared carrier and O a family of functions on A. Write

\[
x\sim_O y \quad\Longleftrightarrow\quad
   o(x)=o(y)\text{ for every }o\in O.
\]

Equality here is metatheoretic equality in each observation's codomain.
If O is empty, every pair is equivalent. This does not imply that the
inscription strings are equal or that the regime contains an equality
operator.

\textbf{Proposition 1 --- separation cannot be recovered from a
collapsed quotient.} Suppose x and y have the same O-profile, while a
task f requires f(x) different from f(y). No function on the current
observational quotient can answer f on all inputs.

\textbf{Proof.} The quotient map q sends x and y to the same class. If
f=g composed with q, then f(x)=g(q(x))=g(q(y))=f(y), a contradiction.
Therefore a representation supporting this task must preserve more
information than q preserves. This may be achieved by adding an
observation that separates the pair. The proposition does not select
that observation. ∎

For \texttt{\textless{}a\textgreater{}} and
\texttt{\textless{}ab\textgreater{}}, a predicate for presence of
\nolinkurl{b} and a predicate for payload length greater than one both
separate the pair, under a supplied string decoder. They disagree on
\texttt{\textless{}aa\textgreater{}}. Full inscription identity also
separates the pair and retains additional information. These are
alternative admitted solutions, not capabilities obtained without a
decoder or comparison operation.

\subsection{5. A history projection can destroy the
answer}\label{a-history-projection-can-destroy-the-answer}

Let \(H\) be a set of admissible histories, \(\sigma:H\to S\) the state
or protected answer required by a task, and \(\pi:H\to V\) the report
shown to an auditor.

\textbf{Theorem 2 --- history report criterion.} An answer map
\(A:\pi(H)\to S\) satisfying \(A(\pi(h))=\sigma(h)\) for every
\(h\in H\) exists if and only if

\[
\pi(h)=\pi(h')\quad\Longrightarrow\quad\sigma(h)=\sigma(h').
\]

When it exists, it is unique on the report image. In particular, two
histories with the same report and different required answers rule out
every universally correct auditor seeing only that report.

\textbf{Proof.} A factorization gives
\(\sigma(h)=A(\pi(h))=A(\pi(h'))=\sigma(h')\). Conversely, define
\(A(\pi(h))=\sigma(h)\); the displayed condition makes this
well-defined. Every member of \(\pi(H)\) has this form, so the factor is
unique. For a finite listed history population with decidable report
equality, the construction is executable: find the first matching report
and return that history's answer. Extending the map outside \(\pi(H)\)
requires a chosen answer there; no such extension is needed for
reconstruction on admissible inputs. ∎

The source establishes the finite natural-number version as
\nolinkurl{Compression.factors_iff_fiberConstant} and
\nolinkurl{factor_unique_on_image}
(\href{lean4/Derivation/Compression.lean\#L167}{source}). The
carrier-general statement and uniqueness on the report image are
mechanized separately as
\nolinkurl{GenesisPaper.history_report_criterion} and
\nolinkurl{history_factor_unique} in the bundled
\nolinkurl{lean4/Paper.lean}. The existence direction uses
\nolinkurl{Classical.choice}; the finite source specialization is
constructive.

This is an information-loss theorem, not a theorem that every task
requires every raw token of history. A sufficient state summary may
preserve \(\sigma\). Conversely, more context alone does not guarantee
that an auditor will use it correctly. Raw history retention is a useful
provenance policy; task-sufficient reconstruction is the mathematical
requirement.

The local response in the trials could be consistent with its asserted
premise and still be inappropriate for the adopted state. Checking
internal steps and checking the binding of their premises to the actual
history are separate tasks. The report criterion makes this distinction
precise for any specified answer function on admissible histories.

\subsection{6. Genesis as an explicit finite admission
study}\label{genesis-as-an-explicit-finite-admission-study}

The \href{docs/sources/genesis_education.txt}{Genesis education account}
tracks the external 19L57G1 study. Here the source is pinned to
\nolinkurl{f1d56c010e6ba10a3b9418bd7cc5953c067250f2}. The earlier sealed
tag \nolinkurl{sealed-derivation-v1} (\nolinkurl{3100c4e}) and September 25
review (\nolinkurl{b368f28c}) retain their own scopes. The education
implementation covers a selected family; it is not a proved copy of
every source research branch.

A finite regime is

\[
R=(D,F,I,U,Q).
\]

The admitted data have indices \(D=\{0,\ldots,n-1\}\), with opaque
presentation labels retained in the token-aware model. The indexed form
list \(F\) assigns a positive arity \(a_f\) to each form. The set \(I\)
contains instances \((f,(d_1,\ldots,d_{a_f}))\) over admitted data;
\(U\) contains forms carrying a uniform all-tuples rule; and \(Q\)
contains supplied requirements \((i,j)\) with \(0\le i<j<n\). Each
requirement is an unordered separation obligation \(\{i,j\}\), stored as
\((\min(i,j),\max(i,j))\); the stored order is not an orientation. These
are the declared host structures through which the primitive operations
are specified.

Regime isomorphisms are bijections of datum and form indices preserving
form arities, ordered instance arguments, uniform-rule membership and
these unordered requirements. Mapped requirement endpoints are sorted
before comparison. Opaque presentation labels are erased before this
comparison and do not constrain the bijections. Thus exchanging the
endpoints of a requirement can identify two responses; an equivalence
relation preserving directed requirements would be a different
convention.

\begin{longtable}[]{@{}
  >{\raggedright\arraybackslash}p{(\linewidth - 2\tabcolsep) * \real{0.5000}}
  >{\raggedright\arraybackslash}p{(\linewidth - 2\tabcolsep) * \real{0.5000}}@{}}
\toprule\noalign{}
\begin{minipage}[b]{\linewidth}\raggedright
Event family
\end{minipage} & \begin{minipage}[b]{\linewidth}\raggedright
Declared effect and admission condition
\end{minipage} \\
\midrule\noalign{}
\endhead
\bottomrule\noalign{}
\endlastfoot
Supply & Leave the admitted regime unchanged; the event may still be
retained in the history. \\
AdmitDatum & Append an opaque label unless it is already admitted;
repeated label admission designates the existing datum. \\
AdmitForm & Append a new form with positive arity, even if another form
has the same arity. \\
AdmitInst & Insert an instance only when the form, arity and datum
indices are valid. \\
AdmitRuleAll & Insert a uniform rule only for an admitted form. \\
Impose separation & Insert the canonical pair \((i,j)\) for an unordered
requirement only when \(0\le i<j<n\). \\
\end{longtable}

Instances, rules and requirements are retained as sets. These event
semantics are defined in the
\href{code/derivation/model.py\#L106}{source model}; ill-formed
admissions are refused. Reading L starts with
\texttt{\textless{}\textgreater{}} supplied but unadmitted. Reading F
also admits that initial token. Both use the same deduplicating
admission convention thereafter. Neither reading is selected by
inspecting the glyphs.

Let \(\operatorname{Ctx}(R)\) contain exactly the patterns whose form is
admitted, whose slot count equals its arity, whose fixed slots name
admitted data, and which contain at least one hole. Plugging \(d\)
replaces every hole by that same datum. Then

\[
R\vdash(f,\mathbf d)
\quad\Longleftrightarrow\quad
(f,\mathbf d)\in I\ \text{or}\ f\in U,
\qquad
\operatorname{Prof}_R(d)
=\{C\in\operatorname{Ctx}(R):R\vdash C[d]\}.
\]

Derivability here is evaluated on well-formed instances. Equal profiles
define observational collapse; a context in their symmetric difference
witnesses separation. With no admitted instances, every datum has the
same profile, even if uniform rules are present: each such rule treats
every plugging alike (\nolinkurl{collapsed_of_no_insts},
\href{lean4/Derivation/Inertia.lean\#L37}{source}). Retaining another
datum alone therefore need not increase expressive power.

The finite context grammar gives the exact count

\[
|\operatorname{Ctx}(R)|
=\sum_{f\in\operatorname{dom}(F)}\bigl((n+1)^{a_f}-n^{a_f}\bigr).
\]

For each form there are \(n+1\) choices per slot, including the hole;
subtract the \(n^{a_f}\) patterns without holes. The companion theorem
\nolinkurl{GenesisPaper.contextWords_length} formalizes the corresponding
recursive enumeration count. \nolinkurl{contextWords_mem} characterizes
exactly the well-formed words containing a hole, and
\nolinkurl{contextWords_nodup} establishes that the enumeration has no
repetitions. Form identifiers keep distinct forms disjoint even when
their arities coincide. These companion proofs are separate from the
frozen source. Its explicit grammar partitions words into hole-free
words and words containing a hole. Python and Rust independently
enumerate the declared contexts, including repeated fixed arguments and
repeated holes. It also explains why large arities can make profile
enumeration expensive.

The source establishes two representation results. \nolinkurl{erase_step}
and \nolinkurl{erase_replay} show that token-aware execution commutes with
erasure to datum counts
(\href{lean4/Derivation/Erasure.lean\#L61}{source}). Erasing a history
is state-dependent: a duplicate datum admission becomes a supply, which
leaves the erased regime unchanged. \nolinkurl{RegimeExt} then identifies
list representations with equal datum counts and form tables and
membership-equal instance, rule and requirement collections;
\nolinkurl{derivable_ext} through \nolinkurl{validResponse_ext} transport
the stated observables (\href{lean4/Derivation/Ext.lean\#L37}{source}).
These are Lean-to-Lean theorems. Executable comparisons with Python and
Rust have their own input populations; a general refinement theorem over
either implementation's source has not been established.

\subsection{7. What ``minimal'' must
include}\label{what-minimal-must-include}

A response is an ordered finite list of \nolinkurl{NewDatum},
\nolinkurl{NewForm}, \nolinkurl{NewInst} and \nolinkurl{NewRuleAll} items. An
instance can refer to a form added earlier in the list. A valid response
applies successfully, separates its target, and leaves all previously
imposed requirements separated. The declared costs are

\[
\mu_1(\Delta)=|\Delta|,\qquad
\mu_2(\Delta)=\left(|\Delta|,
  \sum_{\operatorname{NewForm}(a)\text{ in }\Delta}a\right),
\]

with lexicographic comparison for \(\mu_2\). The sum counts new-form
items, including different forms of equal arity.

The two challenge inputs are explicitly

\[
\begin{aligned}
R_L&:\quad D=\{0,1\},\quad F=I=U=\varnothing,
       \quad Q=\{\{0,1\}\},\\
R_F&:\quad D=\{0,1,2\},\quad F=I=U=\varnothing,
       \quad Q=\{\{1,2\}\}.
\end{aligned}
\]

Datum \(0\) in \(R_F\) is the admitted initial token. The required pairs
are stored as \((0,1)\) and \((1,2)\), respectively. In each input the
target requirement is initially unsatisfied; a valid response must
separate its endpoints.

\textbf{Source result --- exact first-separation cost.} In the initial
form-less, instance-less challenge regimes, the minimum \(\mu_1\) cost
is exactly two in both readings, over the entire declared item grammar
and every positive arity.

\textbf{Proof.} The empty response leaves the pair collapsed. A single
\nolinkurl{NewDatum} or \nolinkurl{NewForm} leaves the instance set empty. A
single \nolinkurl{NewInst} or \nolinkurl{NewRuleAll} cannot apply without an
admitted form. Thus no response of at most one item separates the
target. A fresh unary form followed by an instance at one demanded
endpoint does separate it, and supplies the upper bound. The source
formalizes these steps in \nolinkurl{cost1_insufficient_general},
\nolinkurl{deltaU0_valid} and \nolinkurl{deltaU0F_valid}
(\href{lean4/Derivation/Minimality.lean\#L102}{source}). ∎

\textbf{Source result --- ties and refined uniqueness.} Reading L has
infinitely many nonisomorphic minimum responses under \(\mu_1\). For
each \(k\ge1\), add one \(k\)-ary form and its instance
\((1,\ldots,1)\). The all-hole context separates the two required data.
Distinct arities cannot be identified by a form-arity preserving
isomorphism (\nolinkurl{deltaK_valid}, \nolinkurl{markK_pairwise_noniso}).
The bounded arity-at-most-two census is a finite slice: L has six valid
deltas in three classes; F has ten in five classes
(\href{docs/sources/genesis_source.txt\#L308}{source table}).

Under \(\mu_2\), both readings have the unique minimum \((2,1)\), up to
the source's datum/form isomorphism, which preserves the unordered
requirement. A successful two-item response with arity sum at most one
must add a unary form and one instance. In L either endpoint choice is
isomorphic to the other. In F a mark on the initial, unrequired datum
separates neither demanded endpoint; the two successful endpoint choices
are again isomorphic. These grammar-wide statements are
\nolinkurl{mu2_unique} and \nolinkurl{mu2_unique_F}
(\href{lean4/Derivation/Minimality.lean\#L102}{L proof},
\href{lean4/Derivation/MinimalityF.lean\#L86}{F proof}).

The cost belongs to this accounting language. Bundling a form and
instance into one new primitive would change it. Two admitted items do
not mean two observers, two certificates or two new data at every stage.

\subsection{8. Closure and the next
demand}\label{closure-and-the-next-demand}

\textbf{Source result --- finite round bound.} Fix \(n\) data and a
duplicate-free, well-formed initial requirement store \(Q_0\). Suppose
every round imposes a fresh canonically stored unordered pair and
applies a response containing no \nolinkurl{NewDatum}. Then a run of \(m\)
successful rounds satisfies

\[
2\bigl(|Q_0|+m\bigr)\le n(n-1).
\]

\textbf{Proof.} Responses leave the requirement store unchanged after
imposition. Each fresh imposition therefore adds exactly one pair; the
datum count stays fixed. The final store contains \(|Q_0|+m\) distinct
members of the \(n(n-1)/2\) possible pairs. This is
\nolinkurl{closedRun_round_bound}; from an empty store it gives
\(2m\le n(n-1)\) (\nolinkurl{closedRun_round_bound_from_empty},
\href{lean4/Derivation/Rounds.lean\#L452}{source}). ∎

The declared closure policy chooses a currently collapsed pair and
requires its separation while preserving all prior requirements. In any
well-formed regime satisfying those requirements, a fresh unary form
with one instance at an endpoint supplies such a response. Old-form
derivability remains unchanged, so old witnesses survive
(\nolinkurl{derivable_extend}, \nolinkurl{mark_progress},
\href{lean4/Derivation/Progress.lean\#L43}{source}). Consequently an
unguarded maximal run of this policy ends with all pairs separated: it
cannot continue indefinitely by the bound, and it cannot stop on an
unresolved pair for lack of a response. The companion theorem
\nolinkurl{GenesisPaper.policy_terminal_separates} formalizes this
endpoint implication under well-formedness and preservation of existing
requirements. A bounded implementation can instead stop at a search or
resource limit; that is not the all-pairs-separated result.

Raw admission does not itself enforce this preservation policy. A unary
instance at datum 0 separates data 0 and 1; adding the uniform rule for
that form makes both profiles equal. More admitted structure can
therefore mean fewer observational classes. The round bound tracks
retained requirements, not a monotonicity theorem for arbitrary profile
changes.

The closure policy supplies the demand to separate a collapsed pair. Its
fixed-carrier endpoint does not by itself supply a fresh datum or a new
goal. The later source continuation studied additional, explicitly
declared policies and obtained further results:

\begin{itemize}
\tightlist
\item
  Both authorized continuation attempts closed and were merged. The
  second evaluated 144 engines from each of the two readings, giving 288
  runs, with two-item and arity-two limits, \texttt{max\_data\ =\ 4},
  fifteen rounds, and a 250-node guard per run. All twelve rows of its
  FORCED census were decided under their stated stage or run readings
  (\href{docs/sources/continuation.txt\#L466}{source census}).
\item
  For the cumulative element count \(e(R)=|D|+|F|+|I|+|U|+|Q|\), a
  successful primitive step never decreases \(e\), and a step changing
  the regime increases it by one. Switch events leave the regime
  unchanged. Thus no allowed history can return to a regime it has left,
  and every cycle has a constant regime. These are all-history results,
  including \nolinkurl{ereplay_elems_lt_of_ne}, \nolinkurl{no_return} and
  \nolinkurl{cycle_regime_const}
  (\href{lean4/SecondRun/Run.lean\#L42}{source proof}). They concern
  cumulative admitted elements, not datum count alone or physical
  irreversibility.
\item
  The run-level two-ness property holds for all 64 responding LIVE
  engines, with its witness on every run path before a bound or guard.
  This does not assert two data at every stage. Other census claims have
  their own countermodels and bounds. In particular, one state reached
  after both five and six rounds shows that elapsed round count is not
  generally a function of the regime. Four engines carry the grammar
  forward with a stake and no stream within the measured bounds.
  Continuation beyond those bounds and a minimum number of declared
  inputs remain open
  (\href{docs/sources/continuation.txt\#L792}{terminal account}).
\end{itemize}

These results preserve both the finite constructive progress and its
exact continuation obligations. They do not require reopening the frozen
experiment.

\subsection{9. An admitted witnessed
difference}\label{an-admitted-witnessed-difference}

Once a relation R on a carrier A is explicitly admitted, suppose

\[
\forall x\in A,\;\neg R(x,x),\qquad R(a,b).
\]

\textbf{Proposition 3 --- endpoint distinction.} Then a differs from b,
so the support of this particular edge is the two-element set \{a,b\}.

\textbf{Proof.} If a=b, the supplied witness R(a,b) becomes R(a,a),
contradicting irreflexivity. ∎

This argument does not force irreflexivity from opaque marks. It states
the consequence of an admitted law and an inhabited relation. It says
nothing about the total carrier size, the number of certificates
witnessing the edge, or whether a reverse edge exists. A two-element
carrier with only the edge a→b has one oriented relation occurrence and
no b→a occurrence.

There is also a precise universal property for the chosen witness. Let
\(D_2\) have carrier \(\{0,1\}\) and only the edge \(0\to1\). Given
\(R(a,b)\), the function \(f(0)=a, f(1)=b\) is the unique
endpoint-preserving map from \(D_2\) that preserves its edge. The proof
checks the two inputs; the supplied witness verifies edge preservation.
Under irreflexivity this map is injective by Proposition 3. The source
formalizes \nolinkurl{D2_initial_pointed} and
\nolinkurl{irrefl_witness_ne}
(\href{lean4/Derivation/PrimDiff.lean\#L120}{source}). This represents
the chosen witness inside the relation, not every edge of an arbitrary
\(R\), and does not require a reverse edge or a composition law.

\subsection{10. From Genesis to the presentation
questions}\label{from-genesis-to-the-presentation-questions}

For a presentation S(n), its active demand register and its maps between
stages, the practical questions now become precise:

\begin{enumerate}
\def\labelenumi{\arabic{enumi}.}
\tightlist
\item
  Which data and operations are admitted at the current stage?
\item
  Which declared task cannot be answered with the current observations?
\item
  Which candidate extensions are permitted, and under which costs?
\item
  Which prior answers, inscriptions or algebraic laws must survive?
\item
  What evidence shows both preservation and resolution of the selected
  problem?
\end{enumerate}

Two source results make the preservation questions concrete.

First, let equivalence relations \(\sim_f\) and \(\sim_c\) on the same
set satisfy \(x\sim_f y\Rightarrow x\sim_c y\). The map

\[
q:X/{\sim_f}\longrightarrow X/{\sim_c},\qquad [x]_f\longmapsto[x]_c
\]

is well-defined and surjective. It is injective exactly when the reverse
implication also holds. If an old identification fails in the finer
relation, two distinct fine classes remain over one coarse class.
\textbf{Proof:} use the same representative for surjectivity; equality
of quotient classes gives both directions of the injectivity criterion.
Successive such maps compose without choosing representatives. These are
\nolinkurl{collapseMap_surjective},
\nolinkurl{collapseMap_injective_iff_reverse_refinement} and
\nolinkurl{collapseMap_compose}
(\href{lean4/Derivation/QuotientRetention.lean\#L59}{source}).

Second, for supplied maps \(f:P\to Q\) and \(g:Q\to P\), there exist
injective maps \(i:P\to C\), \(j:Q\to C\) satisfying \(i=j\circ f\) and
\(j=i\circ g\) if and only if

\[
g\circ f=\operatorname{id}_P,\qquad f\circ g=\operatorname{id}_Q.
\]

\textbf{Proof.} Substitute the two identifying equations and cancel the
injective maps to get the round trips. Conversely take \(C=P\),
\(i=\operatorname{id}\) and \(j=g\); the round trips make \(g\)
injective and give the equations. The source theorem is
\nolinkurl{faithful_descent_iff_round_trip_coherent}
(\href{lean4/Derivation/PresentationDescent.lean\#L129}{source}).
Failure for this pair of maps requires retaining their presentation
distinction; it does not rule out another translation.

Canonicalization has a related but separate contract. If a chosen
representative \(C(R)\) lies in the isomorphism orbit of \(R\) and is
constant on each orbit, then \(C(R)=C(S)\) exactly when \(R\cong S\). An
injective serialization of these representatives gives the same
criterion for keys. The source establishes this abstract
\nolinkurl{key_classification}
(\href{lean4/CanonicalMinimizer.lean\#L100}{source}) and establishes
relabelling preserves regime isomorphism
(\href{lean4/Derivation/IsoCanon.lean\#L8}{source}). Connecting an
optimized implementation to it still requires its actual normalization,
serialization and search hypotheses. List order and duplicate
conventions cannot be discarded merely because the abstract theorem
holds.

The companion \nolinkurl{GenesisPaper.serializedRelabel_sound} in
\nolinkurl{lean4/PaperCanonical.lean} formalizes a concrete soundness
direction: if two well-formed regimes have equal sorted five-field
tuples after datum and form relabeling by supplied bijections, they are
isomorphic. Its proof returns the composed isomorphism and establishes
that collection sorting preserves membership. The companion
\nolinkurl{GenesisPaper.Wire.portableBytes_injective} in
\nolinkurl{lean4/PaperWire.lean} formalizes the length-prefixed field
layout and its unsigned 64-bit big-endian encoding. Every encoded
integer, including a collection length, must be less than \(2^{64}\). A
structural decoder recovers the tuple from its words, and the
fixed-width byte encoding retains those words. Thus equal encoded
relabellings imply regime isomorphism under the stated well-formedness,
range and supplied-bijection hypotheses (\nolinkurl{portableBytes_sound}).

The companions \nolinkurl{lean4/PaperOrbit.lean} and
\nolinkurl{lean4/PaperSearch.lean} close the concrete exhaustive-search
classification question. For well-formed regimes with duplicate-free
instance, rule and requirement collections, an isomorphism constructs
bounded datum and form bijections and preserves the entire set of
serialized relabellings. An insertion-based enumeration is proved to
contain exactly all permutations of each finite index range. The
executable minimum search attains a least structural tuple
(\nolinkurl{canonicalKey_attained}), and its keys are equal exactly when
the input regimes are isomorphic
(\nolinkurl{canonicalKey_classification}). Composing this with byte
injectivity gives \nolinkurl{canonicalBytes_classification} under the
unsigned-word bounds. These theorems supply the enumeration and
minimum-selection arguments; they do not assume orbit invariance or an
orbit-choice oracle. General refinement of the actual Python and Rust
source programs, including their optimized searches, remains a separate
engineering obligation.

These criteria permit a mathematical representation to be checked
against its claimed source structure. Selecting a new operation remains
a declared construction decision, with its preservation obligations
stated.

\subsection{11. Engineering follow-up}\label{engineering-follow-up}

There is now a concrete software connection:
\texttt{omega\_wdp::genesis\_construction} calls the canonical
\nolinkurl{omega_ooutwire_primitives} admission implementation and
compares every accepted prefix with WDP's deduplicating execution model,
including all five regime fields. Starting from a literal supply of
\texttt{\textless{}\textgreater{}}, its explicit incidence
interpretation reads admitted instances as simplices, checks vertices
and faces, and delegates exact rational homology to the existing
mathematics code. The triangle boundary has Betti numbers (1, 1);
admitting its face gives (1, 0, 0). The
\href{docs/sources/incidence_construction.txt}{construction account}
states the interpretation and resource limits. This is an implemented
walk with checked executions, not a general Lean refinement theorem over
Rust or a claim that the incidence interpretation is selected by the
token alone.

PFP is the separate engineering follow-up: candidate, checker, receipt
and authorized use must be bound to their task and history. Its intended
consolidation draws on the existing \texttt{011–NNN} source range.
OEWOLAND, OEWAS, OEEIS and OPCIV have no publication commitment here. No
implementation or disclosure of them is required to make the
mathematical statements above intelligible.

\Needspace{480pt}
\subsection{Mechanism inventory}\label{mechanism-inventory}

\begin{longtable}[]{@{}
  >{\raggedright\arraybackslash}p{(\linewidth - 4\tabcolsep) * \real{0.3333}}
  >{\raggedright\arraybackslash}p{(\linewidth - 4\tabcolsep) * \real{0.3333}}
  >{\raggedright\arraybackslash}p{(\linewidth - 4\tabcolsep) * \real{0.3333}}@{}}
\toprule\noalign{}
\begin{minipage}[b]{\linewidth}\raggedright
Mechanism
\end{minipage} & \begin{minipage}[b]{\linewidth}\raggedright
Mathematical evidence
\end{minipage} & \begin{minipage}[b]{\linewidth}\raggedright
Executable scope
\end{minipage} \\
\midrule\noalign{}
\endhead
\bottomrule\noalign{}
\endlastfoot
Observation and history report descent & Carrier-general companion
proofs and pinned finite compression proofs & Finite observation
profiles in Python and Rust; no inference about trial success rates \\
Admission, erasure and collection extensionality & Pinned Lean
transition and representation theorems & Frozen Python model and Rust
companion; the documented name collision distinguishes them \\
Minimum responses and closure & Grammar-wide cost, nonuniqueness,
uniqueness, preservation and round-bound proofs & Arity-at-most-two
minimum census; explicit finite preservation population \\
Chosen witnessed edge and presentation retention & Pinned pointed-graph,
quotient-retention and round-trip theorems & These are mathematical
structure results; no forced physical interpretation \\
Canonical classification & Concrete exhaustive-search and byte-key
classification for well-formed set regimes & Versioned portable protocol
and native Lean/Python/Rust agreement; general source-code refinement
remains open \\
Continuation & Pinned all-history growth and no-return proofs &
Historical census bounds remain those of the cited source \\
Primitive-to-incidence interpretation & Explicit construction choice &
The previously reported WDP triangle executions remain separate from
this standalone primitive reproduction \\
\end{longtable}

\subsection{12. Sources, provenance and verification
status}\label{sources-provenance-and-verification-status}

The \href{source_provenance.json}{source manifest} records the
historical inputs. The mechanism inventory above distinguishes formal
results, computations and construction choices. The complete
\href{docs/source_coverage.md}{source coverage record} inventories 3,147
tracked files and assigns all 121 research directories a disposition.
Inventory coverage is not an independent verification of every proof in
those files.

The \href{docs/sources/closeout.txt\#L490}{final source closeout}
records the author's September 28 freeze of snapshot
\nolinkurl{90ab34a056d80bb1048d926abcfdfef81e578fe1}, approved with merge
\nolinkurl{50d15a44}. Frozen reporting files retain their earlier ``Not
frozen'' headers; the later approval governs their status. Three
independent verdicts stand:

\Needspace{100pt}
\begin{longtable}[]{@{}ll@{}}
\toprule\noalign{}
Axis & Final source status \\
\midrule\noalign{}
\endhead
\bottomrule\noalign{}
\endlastfoot
Research phase & \nolinkurl{CLOSED_WITH_OPEN_OBSTRUCTIONS} \\
Alpha reference implementation &
\nolinkurl{READY_WITHIN_DECLARED_SCOPE} \\
Original seed completion gate & \nolinkurl{UNSATISFIED} \\
\end{longtable}

The completion gate still records unmet live-injection and comparison
clauses, and an unestablished input-count floor. Companion work does not
retrospectively rewrite that frozen verdict. It can establish
additional, separately identified results.

A source implementation defect also remains part of the record. Its
\nolinkurl{delta_events} function generates a supposedly fresh label
\texttt{w\textless{}n\textgreater{}} from the current datum count. If
the sole admitted token is already \nolinkurl{w1}, a \nolinkurl{NewDatum} item
emits \texttt{AdmitDatum("w1")}; deduplication leaves the state
unchanged (\href{code/derivation/candidates.py\#L114}{source}). This was
reproduced during this review. Correcting the Rust implementation is
companion work; comparisons must identify both their collision-free
agreement scope and cases where the corrected behavior differs.

The standalone companion includes unchanged pinned Python and Lean proof
sources, a self-contained Rust crate, and executable reproduction
commands. Lean 4.33.0 compiles the pinned root, parked modules,
canonical-minimizer research modules, second-run kernel and continuation
proofs, and the companion manuscript corollaries. The compiled census
reports 2,475 non-internal theorem constants, including
compiler-generated theorems. Pinned modules use only \nolinkurl{propext}
and \nolinkurl{Quot.sound}; the carrier-general image-factorization
companion also uses \nolinkurl{Classical.choice}. This population is
neither a count of manuscript claims nor verification of all 121 source
research directories.

The standalone replay reproduces the four minimum-response counts in §7,
checks the context formula in 35 cases, and verifies fresh-mark
preservation for 3,996 targets across 1,436 regimes. Native Rust and
frozen Python agree on every context observation and portable key in
that regime population. The native proved Lean minimum search agrees
with their portable keys on those 1,436 labelled regimes; two additional
cases check mixed arities and unsorted list representations. Its
interface rejects 17 malformed or over-budget inputs. The Rust primitive
suite has 135 passing tests; 55 unrelated projection-study tests are
excluded from that result. No new model trial or full frozen-source
completion-gate run is claimed. Commands and scope are in
\href{README.md}{the README}.

Python's historical refined keys sort colors by \nolinkurl{repr}; Rust's
local keys use structural tuple ordering. They can select different
labelled representatives of the same orbit. These local keys and their
derived fingerprints must not be exchanged as identical values. The
companion provides \nolinkurl{portable_key_v1} in both languages:
minimize the structural serialization over all bijection pairs, then use
the versioned, length-prefixed unsigned-integer encoding in
\href{docs/portable_keys.md}{the protocol specification}. Exhausting its
declared permutation budget returns an error. Existing local keys retain
their original conventions. The finite agreement above validates the
tested implementations; it does not assert a compiler-verified
refinement proof for all source code. Both decoders also agree on 72
explicit boundary and malformed-input cases. The compiled Lean wire
encoder agrees with Python on 1,439 tuples: the 1,436 regime keys and
three additional boundary vectors. These comparisons include maximum
unsigned 64-bit words. Decoder acceptance checks the tuple's shape,
well-formedness and collection order; it does not certify an orbit
minimum.

The \href{docs/references.md}{source reader} includes the cited
education chapters, three trial transcripts, historical mathematical
accounts and closeout record. The original repositories are private; the
manuscript's citations point to bundled reading copies or executable
sources. Their manifest records original paths, revisions, Git blob
identities and SHA-256 hashes. The proof and Python sources match the
exact Git bytes at \nolinkurl{f1d56c01}; later companion changes do not
silently change their referents. Repository access is unnecessary for
following these citations.

\end{document}
